This section discusses the results and design choices of the Lapsus fall detection system in relation to the problem statement, requirements, and limitations identified throughout the project. 

The discussion is divided into two main parts. First, the design and technological choices are evaluated, including hardware selection, system architecture, and algorithmic decisions, with a focus on how these choices balance accuracy, power consumption, and usability for elderly users. Second, privacy and security considerations are examined, addressing how sensitive data such as location information, authentication credentials, and alarm notifications are handled and what risks and limitations this introduces.

\subsection{Design and technology}
A \textbf{battery} has been permanently soldered to the microcomputer. This battery does not have a particularly long operational lifetime, however it is rechargeable. Despite this, the wristband must still be recharged regularly. Consequently, it has been essential to design the system running on the wristband to be as energy-efficient as possible. For this reason, a threshold detection algorithm was selected, as it is simple in structure and therefore consumes minimal power.

As stated in the requirement specifications, one of the system requirements was that the wristband should maintain sufficient battery life for an entire day, allowing the user to recharge it overnight. However, this requirement could not be fully met with the battery that was chosen for this project. Battery performance therefore represents an area that should be addressed if further development of the product is to be pursued.

Although the requirement specifications state that the battery life must be a minimum of one day, a preferred target is a battery life of at least one week. This would reduce the risk of the user forgetting to recharge the wristband and thereby lower the likelihood of a fall occurring without being detected. Furthermore, a longer battery life would allow the user to wear the wristband while sleeping, ensuring that any incidents occurring during the night would also be detected.

It would also be beneficial to enable wireless charging for the wristband. This would allow the casing for the micro-controller to be fully enclosed, thereby increasing the potential for making the product \textbf{waterproof} and suitable for use while the user is in the shower. Such a feature would enhance the product's usability and cover a broader range of possible scenarios. However, this was not done due to the need for new components, which were not mandatory for a working prototype and were therefore de-prioritized.

In addition, the prototype is not as discreet as desired. An \textbf{off-the-shelf micro-controller} was used, resulting in several components that are not utilized and could therefore be omitted. This would make it possible to reduce the overall size of the product while also creating space to incorporate additional sensors, thereby increasing the system’s functionality.

One possible enhancement would be the integration of a \textbf{heart rate monitor} capable of tracking the user’s heart rhythm and alerting relatives if the heartbeat were to stop or if irregularities were detected. This would enable the detection of a wider range of potential scenarios and further improve user safety.

The requirement specifications also describe the implementation of an \textbf{inactivity alarm} as a desirable feature. Such an alarm would be triggered if the sensor detected prolonged inactivity by the user. However, this functionality is more challenging to integrate, as the distinction between normal inactivity and situations requiring an alarm can vary significantly from person to person. For this reason, it would be beneficial to implement this feature in conjunction with machine learning, which could analyze individual user behavior and detect anomalies based on personalized activity patterns. As this lies outside the scope of the project, the functionality has not been implemented; however, it represents valuable insight for potential future development.

In the prototype, the wristband is connected to the user’s smartphone via \textbf{Bluetooth}. Communication between the user and their relatives is also conducted through the smartphone. This represents a limitation in situations where the user falls and is unable to move. By contrast, integrating a speaker and microphone directly into the wristband would enable the user to communicate directly with their relatives through the device itself, thereby improving usability and safety in emergency situations.

Furthermore, such integration would allow the system to alert the user when a fall has been detected. This could be achieved by emitting an audible signal from the wristband to indicate that a fall event has been registered. If the user has not fallen, this would provide them with the opportunity to access the application, cancel the alarm, and notify their relatives that they are safe.

This would also contribute to making the wristband more autonomous. The product would be further improved if ownership of a smartphone, or proximity to one, were not required. Bluetooth has an effective range of approximately 10-18 meters, which limits the user’s freedom of movement while wearing the wristband, as they must remain aware of the location of their phone. This limitation may result in alarms not being transmitted to relatives, even if a fall is detected or the emergency button is activated. If the wristband were fully autonomous, overall safety would be enhanced and the product would become accessible to a wider range of users. However, this functionality was not implemented, as it would require the user to purchase a separate \gls{sim} card for the wristband, thereby introducing recurring monthly costs associated with its use. Additionally, such a solution would increase the device’s power consumption, since data transmission via cellular networks requires a stronger transmitter compared to Bluetooth-based communication.

This prototype utilizes Bluetooth Low Energy for data transmission. One advantage of \gls{ble} is that it can remain in a low-power sleep mode when no data is being transmitted and only become active when communication is required. In contrast, transmitting data via cellular networks would require the connection to remain continuously active, which would significantly increase the power consumption of the wristband.

By making the wristband autonomous, the need to adapt the accompanying application to multiple mobile \textbf{operating systems} could be avoided. The prototype has been developed exclusively for use by iPhone users, which excludes a portion of the intended target group. This should therefore be considered in any further development of the product, either by making the application available across multiple smartphone platforms or by designing the system to operate entirely independently of smartphones.

The device has been designed to be worn on the user’s wrist. However, it has not been evaluated whether this location is the most suitable for \textbf{accurate fall detection}. In the context of further development, it is therefore a critical factor to examine where on the body sensors should be placed in order to increase the likelihood of correct fall detection.

As described in the report, the optimal location for sensor placement when detecting falls is around the user’s waist. However, this solution is not ideal for elderly individuals living alone, as such an waist belt may be difficult to put on independently. Additionally, it may be inconvenient for some users if it restricts clothing choices.

The wrist was therefore selected as the placement in this project, as it provides easy access for the user in the event that the emergency button needs to be activated. Furthermore, this choice was appropriate from a design perspective, since one of the requirements was that the product should be as discreet as possible. The prototype thus proposes a design resembling a wristwatch without a display. This ensures that the device does not appear as a conventional fall alarm and can be adapted to the user’s personal style, thereby improving usability and acceptance.

\subsection{Privacy and security}
Privacy and security are central considerations in systems designed for elderly users, particularly when sensitive personal data such as location, health-related events, and contact information are involved. The Lapsus system processes several categories of personal data, including real-time \gls{gps} location, fall events, and phone numbers of relatives, all of which require careful handling to avoid misuse or unauthorized access.

One of the central privacy challenges concerns the handling of \textbf{\gls{gps} location data}. The system continuously collects the user’s location in order to ensure that a real-time position can be sent to relatives together with an alarm. While this enables rapid and accurate assistance, the current implementation does not include mechanisms for automatic deletion of historical location data. As a result, all previously collected \gls{gps} locations remain accessible within the application. This represents a potential privacy risk, as long-term storage of location data may reveal movement patterns and daily routines. In a production-ready system, this would necessitate explicit data retention policies, automatic deletion of outdated data, and user must be able to control and delete stored location information.

Another important security consideration relates to \textbf{user authentication}. The system uses phone number-based login combined with a one-time authorization code sent via \gls{sms}. This method was chosen to reduce cognitive load for elderly users, as it eliminates the need to remember passwords, which can often be a barrier to adoption within the target group. However, this approach is not the most secure authentication method, as it is vulnerable to threats such as \gls{sim} swapping or interception of \gls{sms} messages. While the chosen solution represents a conscious trade-off between security and usability, further development of the system should explore more robust authentication mechanisms that preserve ease of use.

The choice of \gls{sms} as the primary channel for sending \textbf{alarms} to relatives introduces both strengths and weaknesses from a security perspective. \gls{sms} was selected due to its high reliability and accessibility, as it does not require relatives to have the application installed or to enable push notifications. Additionally, \gls{sms} messages can be displayed on in-car infotainment systems, increasing the likelihood that alarms are noticed promptly. However, \gls{sms} communication is inherently insecure and susceptible to spoofing. A malicious actor could potentially imitate the system’s alarm messages and send fraudulent alerts to relatives. Since the alarm messages include a link containing the user’s \gls{gps} coordinates, such an attack could manipulate the link to redirect recipients to false or malicious locations. This highlights the need for additional security measures, such as message verification, digital signatures, or alternative secure communication channels, in future iterations of the system.